\subsection{Retrieval interface and measured operational consequences}
\label{sec:actual-api}
The public source permits both single-time and full-range queries
\citep{hpmapi,hpmcode}. We therefore tested the latter as an operational
comparator instead of treating a two-reading budget as a measured tariff.
The first eligible daily-persistence--MLP task from each external system
was fixed before retrieving optional answers. Two repetitions of five
policies produced 30 trials and 60 actual HTTP requests, including six
common endpoint reads. The protocol limited request starts to one per
second, with no retries or cache-busting. All trials succeeded; timings
include the recorded persistent-connection behavior and uncontrolled
public cache/server state.

\begin{table}[htbp]\centering\small
\caption{Actual public-API measurements on three fixed tasks, two
repetitions per policy. Incremental I/O excludes the separately logged
client throttle wait. Bytes are response-body bytes, not wire traffic.
The full-range interface returns seven hourly samples in one request,
including retained endpoints; it is an operational alternative outside
the two-optional-sample benchmark budget.}\label{tab:actual-api}
\begin{tabular}{lrrr}\toprule
Policy & Requests & Mean I/O (s) & Mean body bytes\\\midrule
Exact adaptive & 2 & 0.497 & 104.3\\
Best deterministic worst-width batch & 2 & 0.539 & 104.7\\
Mean-width-optimal batch & 2 & 0.582 & 104.3\\
Greedy adaptive & 2 & 0.551 & 104.3\\
Full hourly range & 1 & 0.331 & 226.7\\\bottomrule
\end{tabular}\end{table}

All five policies attained the precision target on these three traces.
Bulk retrieval used fewer requests and less observed I/O time, while
also reducing terminal width to zero. Selective retrieval reduced body
size by about 122 bytes; this small difference excludes HTTP/TLS overhead
and does not establish an operational saving. A range returning only the
five internal samples could reduce the bulk response further. Including
the imposed client throttle, mean completion times were 1.984 seconds
for exact adaptation and 1.040 seconds for bulk retrieval; these are not
intrinsic provider latencies. The small workload does not support a
population speedup or monetary-cost claim.

An offline checker matches all 60 received payload hashes to separately
cached source responses, checks all 12 adaptive second-query decisions
and all 30 interval containments (maximum discrepancy
$1.74\times10^{-18}$). Each result can be verified without repeating
provider traffic. The appropriate decision for this observed interface
is bulk retrieval when its unrestricted response is admissible. The
two-query separation instead answers the distinct question posed by a
fixed sample budget; it does not override an available cheaper interface.
\FloatBarrier
